\chapter*{Introducción}
\addcontentsline{toc}{chapter}{Introducción}

% ==========================================================================
% BASE: introducción de la versión 2024 (Paper_Inv_Aplicada_II, pp. 1-3).
% AJUSTES:
%  - El objeto ya no es "el efecto de las multas" en abstracto sino la
%    evaluación de un cambio de política concreto (renovación de la Policía de
%    Tránsito, abril 2022) y su fiscalización diferencial en el espacio.
%  - Anticipar el resultado nulo y su lectura (ausencia de evidencia en un
%    contexto de tratamiento débil), no un hallazgo de "las multas aumentan
%    los accidentes".
%  - Mencionar la nueva fuente de víctimas (FGJ) y la unidad colonia-mes.
% ==========================================================================

\noindent
% --- Párrafo 1: relevancia (reutilizable de 2024) ---
La seguridad vial es un tema central de la política pública a nivel global. De
acuerdo con la Organización Mundial de la Salud, cada año fallecen alrededor de
1.3 millones de personas por hechos de tránsito y decenas de millones más
resultan lesionadas, con costos económicos asociados a la atención médica, la
pérdida de productividad y los daños materiales.

% --- Párrafo 2: la fiscalización como herramienta ---
La imposición de multas por infracciones de tránsito —exceso de velocidad,
conducción bajo los efectos del alcohol, no uso del cinturón— es una de las
herramientas más utilizadas para inducir el cumplimiento de las normas y
disuadir los comportamientos de riesgo. La evidencia sobre su efectividad, sin
embargo, es mixta y sensible al contexto (Capítulo~\ref{chap:literatura}).

% --- Párrafo 3: la pregunta de esta tesis ---
Esta tesis evalúa si el aumento de la fiscalización vial asociado a la
\textbf{renovación de la Policía de Tránsito de la Ciudad de México de abril de
2022} redujo las víctimas de hechos de tránsito. La estrategia explota la
\emph{aplicación diferencial en el espacio} de la reforma: se comparan las
colonias expuestas a la infraestructura de fiscalización preexistente con las no
expuestas, antes y después del cambio, con datos a nivel colonia-mes.

% --- Párrafo 4: qué se encuentra (adelanto) ---
% TODO redactar con cuidado. Borrador:
No se encuentra evidencia robusta de un efecto sobre las víctimas. Las víctimas
fallecidas no cambian; el aumento aparente en las víctimas totales no sobrevive
a un modelo de conteo y coincide con un repunte transitorio de las multas, lo
que apunta a un efecto de reporte. La fiscalización diferencial que induce la
reforma es, además, débil y de corta duración.

% --- Párrafo 5: estructura del documento (reutilizable, actualizar refs) ---
El documento se organiza como sigue. El Capítulo~\ref{chap:literatura} revisa la
literatura sobre fiscalización y seguridad vial. El Capítulo~\ref{sec:contexto}
describe el contexto de la política vial en la Ciudad de México y la reforma de
2022. El Capítulo~\ref{chap:datos} presenta las fuentes y las variables. El
Capítulo~\ref{chap:estrategia} define la estrategia empírica. El
Capítulo~\ref{chap:resultados} presenta los resultados y la robustez. El
Capítulo~\ref{chap:conclusiones} concluye.
